\section{Primitive symbol selection and the remaining density count}\label{sec:selected-dimension}

\status{The integral selection, nonzero multiplier, and normal-degree bound are proved here. The quantitative count is reduced to explicit remaining source statements; it is not claimed proved.}

\paragraph{A residue basis gives an actual integral summand.}
Let $O$ be a DVR, $\mathbb f$ its residue field, and $G$ a set of polynomials supported on a fixed finite set $S$ of monomials. Choose $g_1,\ldots,g_d\in G$ whose reductions form a basis of $\operatorname{span}_{\mathbb f}\overline G$. Their coefficient matrix has a $d\times d$ minor $A$ with unit determinant. If $\pi_S$ selects those coefficient positions, then
\[
 R=A^{-1}\pi_S:O[\mathbf X]\longrightarrow O^d,
 \qquad R(g_i)=e_i.
\]
Thus $O[\mathbf X]=\bigoplus_i Og_i\oplus\ker R$. In particular, the selected summand has a basis of \emph{actual generators}, and its rank is exactly the residue-field dimension. This proves \texttt{RestSel}, by a shorter argument than successive division by a uniformizer. It does not estimate that dimension. \textsf{[SD1]}

A finite saturated polynomial lattice reduces injectively modulo the uniformizer: if $\sum c_i b_i=\varpi P$, saturation puts $P$ in the lattice, and basis coordinates force every $c_i$ to be divisible by $\varpi$. Intersections with monomial-support conditions preserve saturation. For segment products, independence also survives a monomial substitution whenever the map
\[
 (f,(m_j)_j)\longmapsto \ell_f+\sum_j E_jm_j
\]
is injective on the allowed supports: it sends distinct tensor-product monomials to distinct monomials. In the source, this injectivity is obtained by successive division in
$n=\ell+P(u+Qv)$, with $0\le\ell<P$ and $0\le u<Q$; a missing old-frame term has $v=0$. The full geometry and saturation identification for the actual segment spaces remain explicit inputs to this application. \textsf{[SD2]}

\paragraph{The determinant multiplier survives reduction.}
For each late letter, take the last $r$ even selected letters preceding it. They occupy distinct positions modulo $n'$. Indeed, every interval of length $n'$ contains the $r$ periodic translates of one frame; if the chosen $r$ predecessors spanned a whole period, those frame letters and the last predecessor would give at least $r+1$ predecessors after the first. The substituted determinant of these rows is nonzero over every field: setting its diagonal variables to $1$ and other entries to $0$ makes it $1$, even after the prescribed positive powers. Their product $\Delta$ is therefore nonzero modulo $\varpi$. \textsf{[SD3]}

\paragraph{A direct normal-degree bound.}
Put $N=|TT|$, $k=r+K$, $\mathcal L=C_{\rm La}N_s^2\log a$, and
$A=\mathcal L a/(pk)$, $B=ap^{-L_1}$. Count the last $K$ coordinates at even positions and the first $r$ at odd positions. Old-frame letters contribute zero: they are even and have only plane coordinates. At a late letter $x$ of segment $j$, Frobenius substitution multiplies the allowed normal order by $p^{\operatorname{dist}(j,x)}$, giving at most
\[
 p^{\operatorname{dist}(j,x)}o(j,x)\le A+10B.
\]
The late letters across one period partition $TT$; hence segment factors contribute at most $N(A+10B)$. Each label exponent is below a relevant power at most $Q$, so labels contribute at most $NkQ$. Products add these degrees, and sums cannot enlarge support. Consequently every selected-span polynomial satisfies the slightly sharper bound
\[
 \deg_{\rm normal}P\le N(A+10B+kQ)
                 \le10N(A+B+kQ).
\]
This proves the degree part of \texttt{Conclusion2}, with $C_{\rm deg}=10$, under its stated frame and size hypotheses. \textsf{[SD4]}

\paragraph{The exact rank-transfer step.}
Let $\overline\Lambda$ be the span of the source's coarse generators and $\overline N$ the span of the reduced actual symbols. If every coarse generator $u$ satisfies $\overline{\Delta u}=\overline g$ for some actual generator $g$, multiplication by $\overline\Delta\ne0$ injects $\overline\Lambda$ into $\overline N$. SD1 then gives an integral summand of rank at least $\dim\overline\Lambda$. The congruence is the separate statement \texttt{GeneratorModuloUniformiser}; proving a nonzero multiplier alone does not prove it. \textsf{[SD5]}

\paragraph{Two elementary pieces of the remaining count.}
For $r\ge2$, $Q>2b$, and any $c,\rho$, the number of vectors $0\le u_i<Q-2b$ with $c+\sum_i u_i\equiv\rho\pmod Q$ is at least
\[
 Q^{r-1}\left(1-\frac{2rb}{Q}\right).
\]
There are $Q^{r-1}$ unrestricted solutions; for each coordinate, its $2b$ forbidden values account for $2bQ^{r-2}$ solutions. Subtract their union. A terminal positive path step of degree $d$ satisfies
$d\le\mu'_1-\mu_k\le2b$ when its new first weight is $b$ and earlier total negative degree is at most $b$. Thus, if $4Yb\le Q$, applying the count to each of $r$ rows gives the factor
\[
 \left[Q^{r-1}\max\!\left(0,1-\frac r{2Y}\right)\right]^r
\]
used in \texttt{RestMon}. The numerical inequality $4Yb\le Q$ remains to be supplied for the actual parameters. \textsf{[SD6]}

If $J\le N$, $r\ge2$, $k\ge5$, $Y\ge1$, and $1-1/(32Y)\le\alpha\le1$, elementary multiplication and Bernoulli's inequality give
\[
 1-\frac{Nr^2k}{Y}\le
 \alpha^J\left[(1-5/Y)(1-1/Y)(1-r/(2Y))\right]^N
\]
whenever the left side is positive. This accounts for the final source coefficient $C_{\rm rk}=1$; establishing the preceding loss bounds is still necessary. \textsf{[SD7]}

\paragraph{The precise next density target.}
The source chooses $C_{\rm La},C_k\ge1088$, $C_{\rm rk}=1$, $C_{\rm deg}=10$ and proves
\[
 \operatorname{rank}S_g\ge q^{-rK}
 \left(1-\frac{Nr^2k}{Y}\right)
 \prod_{t\in TT}\binom{\deg_t+k-1}{k-1},\qquad q=p^{n'}.
\]
To complete its human proof, expand the coarse-generator congruence; the actual segment-product injection and degree matching; the good-path membership, independence and abundance estimate; and the residue-family summation and numerical bounds. SD1, SD3, and SD4 remove the integral-selection, multiplier-nonvanishing, and degree obligations. SD2, SD5--SD7 give proved intermediate deductions, with their remaining inputs recorded individually. \textsf{[SD8]}
